
\subsubsection{6.1 Why the Proxy Failed}

Canonical HumanEval solutions are carefully crafted reference implementations. They pass all tests by design and follow Python conventions. LLM-generated code---especially early attempts in a repair loop---typically contains more issues: undefined variables, type mismatches, incomplete logic.

Using canonical solutions as an LLM proxy conflates the reference answer with the model's attempts. The 9.15\% warning rate reflects human code quality, not LLM generation behavior.

\subsubsection{6.2 What the Baseline Establishes}

Despite the proxy limitation, the results provide value:

\begin{enumerate}
\item \textbf{Floor established:} Canonical solutions have approximately 9\% warning rate. LLM code is expected to exceed this, providing room for the static analysis intervention to operate.
\end{enumerate}

\begin{enumerate}
\item \textbf{Pipeline validated:} The pylint analysis infrastructure functions correctly. Components are reusable for future work with LLM API access.
\end{enumerate}

\begin{enumerate}
\item \textbf{Warning types identified:} The 9 warning categories detected are actionable for code repair. They provide mechanistic explanations of why code is problematic.
\end{enumerate}

\subsubsection{6.3 Gate Design Implications}

The MUST\_WORK gate served its intended purpose. Without the existence check, the study would have:
\begin{itemize}
\item Expended compute on the wrong code source
\item Produced results that do not test the hypothesis
\item Potentially drawn incorrect conclusions
\end{itemize}

Gate-based validation catches methodology issues early.

\subsubsection{6.4 Limitations}

\textbf{Primary limitation:} No LLM-generated code was tested. API access was unavailable during this study.

\textbf{Secondary limitations:}
\begin{itemize}
\item Single dataset (HumanEval only, not MBPP)
\item Single static analyzer (pylint, not mypy)
\item Single severity filter configuration
\end{itemize}

This is a methodology study establishing baseline and validating infrastructure---not a full hypothesis test.

\subsubsection{6.5 Path Forward}

To complete the hypothesis evaluation:

\begin{enumerate}
\item Re-run h-e1 with LLM API access (GPT-4 or Claude) and measure warning rates on generated code
\item If h-e1 passes (warning rate $\geq 30$\%), proceed to h-m1 (LLM interprets warnings) and h-m2 (non-overlapping errors)
\item Run full evaluation comparing pass@k between static+execution and execution-only conditions
\end{enumerate}

The pipeline, baseline, and methodology are ready. The missing component is LLM API access.
